\section{Motivation}
\label{sec:loosening-motivation}

The translation from B to SMT-LIB developed in \Cref{ch:beer} rests on a tension between the two languages. B is a set-theoretic formalism: as recalled in \Cref{ch:zflean}, every object is ultimately a set, and the distinction between a relation, a partial function and a total function is not one of \emph{kind} but a mere property of one and the same set of pairs. In B, a function \emph{is} its graph. SMT-LIB, by contrast, is many-sorted: every term carries a fixed type, and the translation of B types in SMT types must commit each B construct to one such type. The encoder therefore has to pick, for every B object, a concrete representation---and here lies the difficulty, since a single B object admits several equally faithful SMT representations, between which B operators move freely. Reconciling the rigid type discipline of SMT with the set-theoretic semantics of B is the task of this chapter, and the origin of its title.

Partial functions are the paradigmatic case. Take a B partial function $f : \mathbb{Z} \pfunB \mathbb{Z}$. It is most naturally encoded \emph{functionally}, as a term of SMT type $\intS \toS \optS\,\intS$ sending each integer either to $\someS$ of its image or to $\noneS$. But the same $f$, viewed as its graph, may just as well be encoded \emph{relationally}, as the characteristic predicate of a set of pairs, of SMT type $(\intS \timesS \intS) \toS \boolS$. Both encodings should be denoted by the very same ZFC object---a set of pairs---yet they inhabit different SMT types.

The mismatch surfaces as soon as the two representations meet. Consider the union $f \cupB g$ of two partial functions \(f\) and \(g\) in
\(\mathbb{Z} \pfunB \mathbb{Z}\). Set-theoretically this is the union of their graphs, which need not be deterministic and is in general only a relation. To encode $f \cupB g$ the encoder cannot keep $f$ and $g$ in their functional form $\intS \toS \optS\,\intS$: it must first present each of them relationally, at the common type $(\intS \timesS \intS) \toS \boolS$, and only then take the (relational) union. The same phenomenon governs equality between a function and a relation, membership of a maplet in a set of pairs, and the application of a relation that happens to be a function: each combinator must reconcile the SMT types of its operands before it can act.

Reconciling two types here means moving from a more specific representation to a more general one---from a functional encoding to the relational encoding of the same graph---\emph{without changing the denoted set}. We call this operation \emph{loosening}: the target type is looser in that it forgets the structural guarantees, such as totality or determinism, that the source enjoyed. One may picture it as a kind of \emph{set-level coercion} between SMT representations---a helpful image, even if, strictly speaking, nothing is coerced at the level of sets, where a function already \emph{is} a relation. Loosening is inherently directional: a function may always be viewed as its graph, but a relation is a function only when it is deterministic, so one may always loosen, but never tighten without additional information.

We prove that loosening is \emph{sound} with respect to the ZFC semantics: a loosened term is semantically equivalent to its source, denoting the very same set. This is what licenses the encoder to replace a term by its loosened view whenever two operands must be brought to a common type.

The chapter is organized around the partial order $\sqsubseteq$ that records exactly which types may be loosened into which. Type loosening is the first concrete contribution of the B-to-SMT encoding developed in this thesis: it is the mechanism on which the encoder of \Cref{ch:beer} rests, invoked by every combinator that must reconcile the types of its operands. The development proceeds in three stages. \Cref{sec:loosening-defs} introduces the castable relation $\castable{}{}$ and its constructive companion, the cast paths $\castpath{}{}$, and establishes the order-theoretic structure they induce on $\mathsf{SMTType}$: every type lies in an interval bounded by a \emph{least} and a \emph{most general form}, the two being related by a Galois connection. \Cref{sec:loosening-rules} defines $\mathsf{loosen}$ by recursion on a cast path. \Cref{sec:loosening-correctness} then proves, rule by rule and mirroring the mechanized development, that the generated specification is sound and definite with respect to the semantic cast.
